\newglossaryentry{Labbyn}{
    name=Labbyn,
    description={System do zarządzania infrastrukturą sprzętową tworzony w ramach niniejszej pracy}
}
\newglossaryentry{Prometheus}{
    name=Prometheus,
    description={System telemetrii i monitorowania urządzeń}
}
\newglossaryentry{laboratorium}{
    name=laboratorium,
    description={Pomieszczenie przeznaczone do przechowywania sprzętu komputerowego}
}
\newglossaryentry{rack}{
    name=rack,
    description={Szafa rack, szafa teleinformatyczna, szafa serwerowa, szafa krosownicza}
}
\newglossaryentry{platforma}{
    name=platforma,
    description={sprzęt komputerowy, jednostka centralna}
}
\newglossaryentry{PDU}{
    name=PDU,
    description={Power Distribution Unit, urządzenie wyposażone w wiele gniazd elektrycznych}
}
\newglossaryentry{Ansible}{
    name=Ansible,
    description={narzędzie służące do automatycznego konfigurowania systemów}
}
\newglossaryentry{Playbook}{
    name=Playbook,
    description={Plik w formacie YAML, który definiuje serię zautomatyzowanych zadań do wykonania}
}
\newglossaryentry{open-source}{
    name=open-source,
    description={Rodzaj oprogramowania komputerowego, którego kod źródłowy jest publicznie dostępny}
}
\newglossaryentry{meeting minutes}{
    name=meeting minutes,
    description={Zwięzła notatka służbowa, która dokumentuje najważniejsze ustalenia, decyzje oraz przypisane zadania z przeprowadzonych rozmów}
}
\newglossaryentry{sprint}{
    name=sprint,
    description={Krótka, ograniczona czasowo iteracja, podczas której zespół pracuje nad stworzeniem gotowego, działającego przyrostu produktu}
}
\newglossaryentry{Epik}{
    name=Epik,
    description={Duży zbiór zadań lub funkcja systemu, która jest zbyt obszerna, aby zrealizować ją w jednym sprincie}
}
\newglossaryentry{onboarding}{
    name=onboarding,
    description={Ustrukturyzowany proces wdrażania nowego pracownika do firmy}
}
\newglossaryentry{POC}{
    name=POC,
    description={(ang. Proof of Concept) prototyp potwierdzający, że dany pomysł jest możliwy do wykonania i sensowny z biznesowego punktu widzenia}
}
\newglossaryentry{Pull Request}{
    name=Pull Request,
    description={Mechanizm który powiadamia o zakończeniu prac nad funkcją i prosi o przejrzenie oraz scalenie zmian do docelowej gałęzi repozytorium}
}
\newglossaryentry{SQL Alchemy}{
    name=SQL Alchemy,
    description={Biblioteka języka Python umożliwiająca komunikację z bazami danych oraz mapowanie obiektowo-relacyjne}
}
\newglossaryentry{Backend}{
    name=Backend,
    description={“Serce” systemu odpowiedzialne za logikę, obliczenia, bezpieczeństwo i dostarczanie danych które mają zostać odebrane oraz przekazane użytkownikowi}
}
\newglossaryentry{API}{
    name=API,
    description={Zestaw reguł pozwalający na komunikację pomiędzy poszczególnymi częściami aplikacji}
}
\newglossaryentry{FastAPI}{
    name=FastAPI,
    description={Framework języka Python, służący do budowania zaawansowanych interfejsów API}
}
\newglossaryentry{Pydantic}{
    name=Pydantic,
    description={biblioteka Python, odpowiedzialna za walidację danych wejściowych do backendu}
}
\newglossaryentry{OpenAPI}{
    name=OpenAPI,
    description={ujednolicony standard interfejsów REST API}
}
\newglossaryentry{ORM}{
    name=ORM,
    description={(ang. Object-Relational Mapping) technika mapowania obiektów w kodzie na tabele w bazie}
}
\newglossaryentry{PostgreSQL}{
    name=PostgreSQL,
    description={obiektowo-relacyjny system zarządzania bazą danych}
}
\newglossaryentry{Redis}{
    name=Redis,
    description={baza danych przechowująca dane w pamięci RAM}
}
\newglossaryentry{race condition}{
    name=race condition,
    description={błąd logiczny, w którym wynik operacji zależy od nieprzewidywalnej kolejności wykonania współbieżnych procesów}
}
\newglossaryentry{CORS}{
    name=CORS,
    description={(z ang. Cross-Origin Resource Sharing), mechanizm bezpieczeństwa przeglądarek, który pozwala na określenie konkretnych reguł na podstawie których można odbierać żądania użytkownika}
}
\newglossaryentry{Router}{
    name=Router,
    description={Komponent FastAPI, który pozwala na tworzenie interfejsów API z mniejszych, niezależnych modułów}
}
\newglossaryentry{Pliki statyczne}{
    name=Pliki statyczne,
    description={pliki które nie ulegają zmianom w trakcie działania aplikacji}
}
\newglossaryentry{Zmienne środowiskowe}{
    name=Zmienne środowiskowe,
    description={globalne wartości konfiguracyjne, które pozwalają na zmianę działania systemu}
}
\newglossaryentry{atrybuty obliczalne}{
    name=atrybuty obliczalne,
    description={metoda klasy, która może być wywoływana jako zwykła zmienna}
}
\newglossaryentry{Blokada optymistyczna}{
    name=Blokada optymistyczna,
    description={Strategia kontroli, która sprawdza numer wersji przy próbie zapisu, by upewnić się nie dane nie zostały zmodyfikowane w międzyczasie}
}
\newglossaryentry{Grafana}{
    name=Grafana,
    description={platforma służaca do wizualizacji, monitorowania i analizy danych w czasie rzeczywistym}
}
\newglossaryentry{MVP}{
    name=MVP,
    description={(z ang. Minimum Viable Product) najprostsza wersja nowego produktu lub usługi, posiadająca tylko niezbędne funkcje}
}
\newglossaryentry{RBAC}{
    name=RBAC,
    description={(z ang. Role-Based Access Control) model zarządzania uprawnieniami, które są przyznawane na podstawie ról}
}
\newglossaryentry{Architektura wielotenantowa}{
    name=Architektura wielotenantowa,
    description={architektura, w której różne grupy użytkowników działają w ramach jednej bazy i instancji aplikacji}
}
\newglossaryentry{Alembic}{
    name=Alembic,
    description={narzędzie do migracji bazy danych dla \gls{SQL Alchemy}. Pozwala na wersjonowanie struktur tabel}
}
\newglossaryentry{PromQL}{
    name=PromQL,
    description={(z ang. Prometheus Quesry Language) język zapytań wykorzystywany przez serwer \gls{Prometheus}}
}
\newglossaryentry{Cachowanie}{
    name=Cachowanie,
    description={technika tymaczasowego przechowywania danych często używanych}
}
\newglossaryentry{JSON}{
    name=JSON,
    description={(z ang. JavaScript Object Notation) format danych używany jako standard komunikacji w \gls{API}}
}
\newglossaryentry{JSONB}{
    name=JSONB,
    description={(z ang. \gls{JSON} Binary) typ danych w bazie \gls{PostgreSQL}, który przechowuje format \gls{JSON} w formie binarnej}
}
\newglossaryentry{Websocket}{
    name=Websocket,
    description={protokół komunikacyjny, zapweniający dwukierunkowy transfer danych, przez pojedyncze połączenie w czasie rzeczywistym}
}
\newglossaryentry{CSV}{
    name=CSV,
    description={(z ang. Comma Separated Values) format pliku tekstowego do przechowywania danych, w których wartości są kolejno rozdzielone separatorem}
}
\newglossaryentry{URL}{
    name=URL,
    description={(z ang. Uniform Resource Locator) ujednolicony format adresowania w Internecie, wskazujący dokładną lokalizację zasobu}
}
\newglossaryentry{Endpoint}{
    name=Endpoint,
    description={adres \gls{URL} w \gls{API}, pod którym serwer wykonuje okreśnole funkcje}
}
\newglossaryentry{SQL Injection}{
    name=SQL Injection,
    description={technika ataku polegająca na wstrzykiwaniu złośliwego kodu SQL}
}
\newglossaryentry{CI/CD}{
    name=CI/CD,
    description={metodologia automatycznego testowania kodu, sprawdzania spójności, generowania plików i/lub aktualizacji systemu automatycznie}
}
\newglossaryentry{pytest-html}{
    name=pytest-html,
    description={wtyczka (plugin) do frameworka pytest w języku Python, która umożliwia generowanie raportów z testów w formacie HTML}
}
\newglossaryentry{Selenium}{
    name=Selenium,
    description={framework służący do automatyzacji przeglądarek internetowych}
}
\newglossaryentry{Pytest}{
    name=Pytest,
    description={framework do testowania kodu w języku Python}
}
\newglossaryentry{Selenium WebDriver}{
    name=Selenium WebDriver,
    description={narzędzie i interfejs API służący do auomtayzacji przeglądarek internetowych}
}
\newglossaryentry{WebDriverWait}{
    name=WebDriverWait,
    description={klasa framework'a Selenium służąca do implementacji jawnego czekania, pozwalająca na zatrzymanie wykonywania kodu i poczekaniu na spełnienie określonego warunku}
}
\newglossaryentry{E2E}{
    name=E2E,
    description={(z ang. End-To-End) metoda testowania oprogramowania, która sprawdza działanie całej aplikacji od poczatku do końca, symulując rzeczywiste przypadki użycia}
}
\newglossaryentry{Coding Guidelines}{
    name=Coding Guidelines,
    description={zbiór zasad, konewncji oraz najlepszych praktyk, których programiści przestrzegają podczas pisiania kodu źródłowego}
}
\newglossaryentry{docstring}{
    name=docstring,
    description={opis w kodzie źródłowym, który służy do dokumentowania konkretnego segmentu kodu}
}
\newglossaryentry{PEP 8}{
    name=PEP 8,
    description={oficjalny przewodnik po stylu w języku Python}
}
\newglossaryentry{Docker}{
    name=Docker,
    description={narzędzie do konteneryzacji, pozwala na pakowanie aplikacji z jej zależnościami w lekkie przenośne kontenery}
}
\newglossaryentry{kontener}{
    name=kontener,
    description={lekka i przenośnia instancja oprogramowania odizolowana od systemu}
}
\newglossaryentry{manifest}{
    name=manifest,
    description={plik konfiguracyjny zawierające kluczowe informacje o danym obiekcie lub aplikacji}
}
\newglossaryentry{skrypt}{
    name=skrypt,
    description={zbiór instrukcji stworzony do zautomatyzowania konkretnej czynności}
}
\newglossaryentry{repozytorium}{
    name=repozytorium,
    description={centralne miejsce do przechowywania, zarządzania i udostępniania danych, plików lub kodu}
}
\newglossaryentry{system kontroli wersji}{
    name=system kontroli wersji,
    description={system pozwolający na kontrolowanie pracy nad projektem oraz śledzenie historii zmian projektu}
}
\newglossaryentry{merge}{
    name=merge,
    description={dołączenie zmian kodu z jednej gałęzi do innej}
}
\newglossaryentry{Git}{
    name=Git,
    description={najpopularniejszy system kontroli wersji wykorzystywany na szeroką skalę zarówno przy małych, domowych projektach jak również w dużych organizacjach}
}
\newglossaryentry{GitHub}{
    name=GitHub,
    description={platforma pozwalająca przechowywać repozytorium w chmurze. Działa w oparciu o system Git. Umożliwia wspólną pracę zespołów z dowolnego miejsca}
}
\newglossaryentry{GitHub Actions}{
    name=GitHub Actions,
    description={wbudowana w platformę GitHub obsługa CI/CD, która umożliwia automatyzację przepływów pracy bezpośrednio w repozytorium}
}
\newglossaryentry{runner}{
    name=runner,
    description={maszyna wykorzystywana do wykonywania workflow'ów i przetwarzania wyników}
}
\newglossaryentry{workflow}{
    name=workflow,
    description={zbiór zadań wykonujących instrukcje krok po kroku w celu zautomatyzowania danej czynności}
}
\newglossaryentry{job}{
    name=job,
    description={zestaw kroków wykonywanych w ramach jednego workflow na jednym runnerze np. pobierz kod, zbuduj aplikację, przetestuj aplikację, wyślij raport}
}
\newglossaryentry{step}{
    name=step,
    description={najmniejsza jednostka wykonywalna w obrębie job'a, odpowiada za zrobienie jednej czynności np. pobierz kod albo uruchom skrypt}
}
\newglossaryentry{React}{
    name=React,
    description={biblioteka języka JavaScript służaca do tworzenia intefrejsów użytkownika}
}
\newglossaryentry{Vagrant}{
    name=Vagrant,
    description={narzędzie umożliwiające łatwe tworzenie i zarządzanie maszynami wirtualnymi}
}
\newglossaryentry{IaC}{
    name=IaC,
    description={(z ang. Infrastructure as Code) praktyka polegająca na zarządzaniu i tworzeniu zasobów komputerowych w oparciu o pliki konfiguracyjne trzymane w czytelnych dla człowieka plikach (Manifest)}
}
\newglossaryentry{hot-reload}{
    name=hot-reload,
    description={natychmiastowe wprowadzenie i wyświetlanie zmian w programie bez konieczności zatrzymywania i ponownego budowania}
}
\newglossaryentry{framework}{
    name=framework,
    description={gotowy zestaw narzędzi, bibliotek i reguł, który służy do szybszego tworzenia aplikacji}
}
\newglossaryentry{N-Tier}{
    name=N-Tier,
    description={model projektowania oprogramowania typu klient-serwer, w którym aplikacja zostaje podzielona na niezależne warstwy.}
}
\newglossaryentry{Docker Compose}{
    name=Docker Compose,
    description={narzędzie do definiowania i uruchamiana wielokontenerowych aplikacji Docker}
}
\newglossaryentry{mockup}{
    name=mockup,
    description={(z ang. makieta/wizualizacja) model prezentujący jak projekt będzie wyglądał w rzeczywistości.}
}
